\documentclass[11pt]{article}
\usepackage[margin=1in]{geometry}
\usepackage{amsmath,amssymb,amsthm}
\usepackage{graphicx}
\usepackage[hidelinks]{hyperref}
\usepackage{xcolor}

\newtheorem{theorem}{Theorem}
\newtheorem{remark}{Remark}
\newcommand{\M}{M}
\newcommand{\C}{\mathbb C}

\title{A Literal Answer to Problem 1 in\\
\emph{Quantum Expanders and Geometry of Operator Spaces}}
\author{Automated mathematics run \texttt{fa\_banach\_001}}
\date{9 August 2026}

\begin{document}
\maketitle

\begin{center}
\fbox{\parbox{0.91\textwidth}{\textbf{Status.}
Candidate full answer to the statement as printed. The likely intended
``for all sufficiently large $N$'' version is not answered.}}
\end{center}

\section{Source and printed problem}

The source is G. Pisier, \emph{Quantum Expanders and Geometry of Operator
Spaces}, arXiv:1209.2059, Section 3, Problem 1, page 25 of the source PDF
\cite{Pisier}. For a finite-dimensional operator space $E$, an integer
$N\geq 1$, and $C>1$, the number $k_E(N,C)$ is the least $k$ for which there
are linear maps $f_j:E\to \M_N$ such that, for every $x\in \M_N(E)$,
\[
 \sup_{1\leq j\leq k}\|({\rm Id}\otimes f_j)(x)\|
 \leq \|x\|
 \leq C\sup_{1\leq j\leq k}\|({\rm Id}\otimes f_j)(x)\|.
\]

The printed question is:

\begin{quote}
Let $C>1$. Assume that a finite dimensional space $E$ satisfies
$k_E(N,C)\leq 1$ for all $N$. What does that imply on $E$? Is $E$ exact
with a control on its exactness constant?
\end{quote}

\begin{figure}[ht]
  \centering
  \includegraphics[width=\textwidth]{figures/open_problem_crop.png}
  \caption{The complete problem list on source page 25; Problem 1 is the
  statement answered here.}
\end{figure}

\section{Proof intuition}

The universal quantifier contains the scalar level $N=1$. At that level a
single map into $\M_1=\C$ must recover the norm of every vector of $E$ up to
the factor $C$. Its lower estimate makes it injective, and an injective linear
map from $E$ into a one-dimensional space forces $E$ itself to have dimension
at most one. One-dimensional operator spaces have no residual matricial
complexity: their matrix norms are exactly the scalar matrix norms, so they
are completely isometric to $\C$.

\section{Formal result}

\begin{theorem}
Let $E$ be a finite-dimensional operator space and let $C>1$. If
\[
 k_E(N,C)\leq 1\qquad\text{for every integer }N\geq 1,
\]
then $\dim E\leq 1$. If $E\neq\{0\}$, then $E$ is completely isometric to
$\C$ and its exactness constant is
\[
 \operatorname{ex}(E)=1.
\]
In particular, the printed Problem 1 has an affirmative answer with optimal
control independent of $C$.
\end{theorem}

\begin{proof}
Set $N=1$. Since $\M_1(E)=E$ and $\M_1=\C$, the hypothesis
$k_E(1,C)\leq 1$ supplies a linear map $f:E\to\C$ such that
\[
 |f(x)|\leq \|x\|_E\leq C|f(x)| \qquad (x\in E).
\]
If $f(x)=0$, the right-hand inequality gives $\|x\|_E=0$, hence $x=0$.
Thus $f$ is injective. Since its range lies in the one-dimensional vector
space $\C$, it follows that $\dim E\leq 1$.

Assume now that $E\neq\{0\}$. Realize $E$ concretely as a subspace of
$B(H)$ and choose $e\in E$ with $\|e\|=1$. The map
\[
 u:\C\longrightarrow E,\qquad u(\lambda)=\lambda e,
\]
is an isometry. At every matrix level $m$ and for every scalar matrix
$a=[\lambda_{pq}]\in\M_m$,
\[
 \|u_m(a)\|_{\M_m(E)}
 =\|[\lambda_{pq}e]\|_{B(\C^m\otimes H)}
 =\|a\otimes e\|
 =\|a\|_{\M_m}\|e\|
 =\|a\|_{\M_m}.
\]
Therefore $u$ is a complete isometry, so $E$ is completely isometric to the
subspace $\C=\M_1$ of a matrix algebra. Hence $\operatorname{ex}(E)\leq 1$.
By the normalization of the exactness constant it is always at least $1$,
and consequently $\operatorname{ex}(E)=1$.
\end{proof}

\section{Scope and formulation caveat}

Immediately before introducing matricially subGaussian spaces, the source
observes that a $C$-exact space has $k_E(N,c)=1$ for all sufficiently large
$N$ whenever $c>C$. This makes it plausible that Problem 1 intended the
eventual condition
\[
 k_E(N,C)\leq 1\quad\text{for all sufficiently large }N,
\]
rather than the universal condition printed on page 25. The proof above uses
$N=1$ essentially and says nothing about that eventual version. The result
should therefore be read as both a complete literal answer and a formulation
correction.

\section{Verification and novelty check}

The proof uses only the source definition and elementary properties of
concrete one-dimensional operator spaces. The source TeX and PDF were checked
to confirm that Problem 1 literally says ``for all $N$'' and that $N$ ranges
over positive integers in the definition.

A bounded search was performed in the run's solution, attempt, proof-gap, and
registry indexes; the local arXiv source corpus; the source's sequel and
related exact-operator-space papers; and targeted web searches for the exact
notation $k_E(N,C)$, matricially subGaussian spaces, and later discussions of
the problem. No explicit published observation of the $N=1$ collapse was
found. Novelty confidence is high for the literal formulation correction but
low for any substantive intended problem, since the eventual-$N$ question is
not resolved here.

Human review should focus on whether an implicit convention in the paper was
meant to exclude $N=1$. If so, the packet remains correct only as a diagnosis
of the printed wording.

\begin{thebibliography}{9}
\bibitem{Pisier}
G. Pisier,
\emph{Quantum Expanders and Geometry of Operator Spaces},
J. Eur. Math. Soc. 16 (2014), 1183--1219;
arXiv:1209.2059.
\end{thebibliography}

\end{document}

